\documentclass[a4paper]{article}

\usepackage[spanish]{babel}
\usepackage{graphicx}
\usepackage{amsmath, amssymb}
\usepackage[margin=2cm]{geometry}
\usepackage{fancyhdr}
\usepackage{enumerate}
\usepackage[shortlabels]{enumitem}
\usepackage{parskip}
\usepackage[most]{tcolorbox}
\usepackage[hidelinks]{hyperref}
\usepackage{float}
\usepackage{booktabs}



% cabecera
\pagestyle{fancy}
\fancyhead[l]{Mecánica Celeste}
\fancyhead[c]{Clase \#4}
\fancyhead[r]{\today}
\fancyfoot[c]{\thepage}
\renewcommand{\headrulewidth}{0.2pt} % linea horizontal


\begin{document}

\section{Solución Numérica al Problema de los $N$-Cuerpos}

Dada la imposibilidad de encontrar una solución analítica práctica para $N > 2$ (debido a la lentitud de convergencia de las series de Sundman o Wang), la mecánica celeste moderna se apoya en la \textbf{solución numérica}. Este enfoque permite construir "laboratorios virtuales" para poner a prueba teorías dinámicas y realizar predicciones precisas en ingeniería aeroespacial.

\section{Linealización y Reducción del Orden}

Para resolver numéricamente el sistema de ecuaciones diferenciales de segundo orden que describe el movimiento gravitacional:
\begin{equation}
    \left\{ \ddot{\vec{r}}_i = -\sum_{j=1, j \neq i}^{N} \frac{\mu_j}{r_{ij}^3} \vec{r}_{ij} \right\}_{i=1...N}
\end{equation}

Es necesario realizar una \textbf{reducción del orden} para convertirlo en un sistema de ecuaciones de primer orden. Siguiendo una idea original de Euler, definimos variables auxiliares de velocidad $\vec{v}_i \equiv \dot{\vec{r}}_i$, lo que transforma el sistema de $3N$ ecuaciones de segundo orden en $6N$ ecuaciones de primer orden:

\begin{equation}
    \begin{cases} 
    \dot{\vec{r}}_i = \vec{v}_i \\
    \dot{\vec{v}}_i = -\sum_{j \neq i} \frac{\mu_j}{r_{ij}^3} \vec{r}_{ij} 
    \end{cases}
\end{equation}

Matemáticamente, esto se gestiona mediante un \textbf{vector de estado} $Y$ que compila todas las componentes de posición y velocidad del sistema:
\begin{equation}
    Y = \{x_1, y_1, z_1, ..., x_N, y_N, z_N, v_{x1}, v_{y1}, v_{z1}, ..., v_{xN}, v_{yN}, v_{zN}\}
\end{equation}

\section{El Problema de Valor Inicial (IVP)}

La integración numérica requiere definir \textbf{condiciones iniciales} precisas en un instante $t_0$. Esto implica conocer $6N$ valores: las posiciones $\{ \vec{r}_i(t_0) \}$ y las velocidades $\{ \dot{\vec{r}}_i(t_0) \}$ de todas las partículas. El objetivo es encontrar el estado del sistema en un tiempo futuro $t_0 + \Delta t$.

\section{Unidades Canónicas y Escalas Gravitacionales}

En el Sistema Internacional, la constante $G \approx 6.67 \times 10^{-11} \, \text{m}^3\text{kg}^{-1}\text{s}^{-2}$ es extremadamente pequeña, mientras que las masas y distancias astronómicas son enormes (p.e. $M_{\odot} \approx 10^{30}$ kg). Para mejorar la estabilidad numérica y evitar errores de redondeo, se utilizan \textbf{unidades canónicas} donde $G = 1$.

Un sistema de unidades $(L, M, T)$ es canónico si sus factores de conversión $(U_L, U_M, U_T)$ satisfacen:
\begin{equation}
    1 = \frac{G U_M U_T^2}{U_L^3} \implies U_T = \sqrt{\frac{U_L^3}{G U_M}}
\end{equation}

En estas unidades, el parámetro gravitacional $\mu_j = G m_j$ se reduce numéricamente a la masa $m_j$.

\section{Algoritmos de Integración}

\subsection{Diferencias Finitas y Método Leapfrog}
Una técnica común para propagar el estado del sistema es el uso de \textbf{diferencias finitas}. Un algoritmo destacado es el \textbf{Método de Leapfrog}, una variante de los métodos de Delambre o Verlet, que actualiza posiciones y velocidades de forma alternada para conservar mejor la energía del sistema:
\begin{equation}
    \begin{cases} 
    \vec{v}_i(t_0 + \Delta t) \approx \vec{v}_i(t_0) + \vec{a}_i(t_0) \Delta t \\
    \vec{r}_i(t_0 + \Delta t) \approx \vec{r}_i(t_0) + \vec{v}_i(t_0 + \Delta t) \Delta t 
    \end{cases}
\end{equation}

\subsection{Implementación Computacional}
En la práctica, se utilizan integradores robustos como \texttt{odeint} (basado en ODEPACK). El proceso algorítmico general implica:
\begin{enumerate}
    \item Definir el sistema de ecuaciones diferenciales reducidas (EDMR).
    \item Especificar el estado inicial $Y_0$ y el intervalo de tiempo.
    \item Extraer las trayectorias y reconstruir los vectores de posición y velocidad.
\end{enumerate}

\section{Validación de la Solución}

Una solución numérica es confiable solo si preserva las propiedades físicas globales del sistema. Se deben verificar las \textbf{constantes de movimiento} (cuadraturas) a lo largo del tiempo:
\begin{itemize}
    \item \textbf{Momentum Lineal:} $\vec{P}_{CM} = \sum m_i \dot{\vec{r}}_i$.
    \item \textbf{Momentum Angular:} $\vec{L} = \sum m_i \vec{r}_i \times \dot{\vec{r}}_i$.
    \item \textbf{Energía Total:} $E = K + U$.
    \item \textbf{Teorema del Virial:} En sistemas ligados, el promedio de la energía cinética debe ser $-\frac{1}{2}$ del promedio de la energía potencial ($\langle K \rangle = -\frac{1}{2} \langle U \rangle$).
\end{itemize}

\bibliographystyle{plainurl}
\bibliography{bibliografia} 
\end{document}
